\section{What the family conclusion now supplies}

% BEGIN J78H1
\paragraph{J78H1. The structural and numerical fields at one threshold.}
Under F78D2, fix $C_2\ge1$. There are a single $C_s\ge1$
and a single $x_0$ for which the same $X$ satisfies nine fields
of the original family conclusion at every qualified $x,i$:
\[
 \mathrm{pos},\quad\mathrm{pass},\quad\mathrm{mult},\quad
 \mathrm{rank},\quad\mathrm{kill},\quad\mathrm{num},\quad
 \mathrm{exc},\quad\mathrm{par},\quad\mathrm{small}.
\]
In addition $X.\mathrm{exc}$ is the specific exceptional set of
the arithmetic parameters for that member.

Choose the uniform parameter constant of the parameter chapter, then
take $C_s$ at least that constant, $3C_2+1$ and $22365$.
Use DFD14 at this fixed $C_s$. Take the maximum of the thresholds for
DFD14, F78D2, the uniform parameter inequalities and the exceptional
set estimates. For $x$ beyond that maximum and $i$ at that scale,
first choose the arithmetic exceptional set from the exceptional-set
construction. Then use it as the argument $S$ of F78D2.
The equality $X.\mathrm{exc}=S$ transports its cardinality and element
bounds to $X$. The parameter result supplies all eight parameter
conditions at $C_s$, and DFD14 gives $\mathrm{small}$ at $C_s$.
F78D2 gives killing and F78D3 gives the five numerical clauses there.

For the other four fields, $1\le C_s$ implies positivity in the
source's sense $0\le C_s$. The second part of
$\mathrm{FactsF}(x,i)$ includes its commuting-endomorphism condition.
Apply the cut-defect inequality LN5 with this endomorphism to get
$\mathrm{adm1}$, the passage field on the original denominator. The
localization uniqueness and good-row identity from the two conclusions
in the member's facts, its isotypic internal decomposition, and its
rank comparison give the multiplicity estimate LN6. Finally
the core part of the very same facts supplies $\mathrm{adm3}(L)$,
the rank field. All the constants were fixed before $x,i$.
Increasing $C_s$ before choosing the DFD14 cutoff is allowed, so
further fixed lower bounds on it can be imposed at this stage.
% END J78H1

% BEGIN J78H2
\paragraph{J78H2. The polynomial estimate needed at the two ramified primes.}
On the hub of F78D2 let $\ell$ equal either of the two primes $q,d$.
Suppose its native distinguished-prime estimate
$\mathrm{Conclusion2}(X.\rho,L,A,Q)$ is supplied, where $Q$ is a
polynomial in $f,r,t$ with nonnegative integer coefficients.
For
\[
 C_s\ \ge\ \max\left\{A,1,\deg Q,\sum_\alpha [X^\alpha]Q\right\}
\]
one obtains the full $\mathrm{StarG}(L,X.\rho,C_s)$ field at both
primes. In particular for the source choices $A=2080$ and
$Q=260 f^2t^2(r+t)$ this size comparison is already implied by
the common numerical row bound in F78D3.

Indeed fix a twisted index at $\ell$ and one of its good points.
The host bounds make $\ell$ an odd prime, and the original parameter
choice gives $x^{2C_t+5}<\ell$. The killing condition on this very
point supplies finiteness of the Selmer term required in
$\mathrm{Conclusion2}$. The odd finiteness hypothesis gives finite
Tate--Shafarevich lengths for every integer $0<e<t$.
Thus the assumed polynomial error estimate applies at the point.
Its Tate--Shafarevich sum ranges over just $1\le e<t$, its class
number term is $v_\ell(h_E)$, and its last correction is the same
$\varepsilon_\ell$ appearing in the target inequality.

Put $u=x^\varepsilon=\log b>0$ and
$B=r(b/\ell+S_c+t+\eta+u)+mb\ge0$ in the table notation.
The part of the bound left to compare is
$AfB+Q(f,r,t)u$. All three ranks are at least one. Write $H=frt$.
Each monomial in the support of $Q$ has value at most $H^{\deg Q}$,
because each variable is at most $H$ and $H\ge1$. Summing with
its coefficients gives $Q(f,r,t)\le C_s H^{C_s}$, including
the zero polynomial. Hence the error is at most
$C_sfB+C_s H^{C_s}u$, precisely $\mathrm{midG}(L,C_s,\ell)$.
Adding back the common terms proves the stated bound for every
guarded point. For the specified polynomial the coefficient sum is
$520$ and the total degree is $5$, so the displayed maximum is
$2080$. This accounts for all numerical enlargements in this
conversion. The arithmetic proof of its distinguished-prime
estimate must still be supplied on the chosen representation.
% END J78H2
